\documentclass[11pt,a4paper]{article}

% -------------------------------------------------
% Packages
% -------------------------------------------------
\usepackage[margin=2.2cm]{geometry}
\usepackage{kotex}
\usepackage{amsmath,amssymb}
\usepackage{booktabs}
\usepackage{array}
\usepackage{xcolor}
\usepackage{listings}
\usepackage[most]{tcolorbox}
\usepackage{enumitem}
\usepackage{hyperref}
\usepackage{fancyhdr}
\usepackage{titlesec}

% -------------------------------------------------
% Colors
% -------------------------------------------------
\definecolor{codebg}{RGB}{247,247,247}
\definecolor{codeframe}{RGB}{210,210,210}
\definecolor{keywordcolor}{RGB}{0,70,160}
\definecolor{commentcolor}{RGB}{0,120,80}
\definecolor{stringcolor}{RGB}{170,50,50}
\definecolor{titleblue}{RGB}{35,75,130}

% -------------------------------------------------
% Hyperlinks
% -------------------------------------------------
\hypersetup{
    colorlinks=true,
    linkcolor=titleblue,
    urlcolor=titleblue
}

% -------------------------------------------------
% Header / Footer
% -------------------------------------------------
\pagestyle{fancy}
\fancyhf{}
\lhead{Python Programming}
\rhead{Day 1}
\cfoot{\thepage}

% -------------------------------------------------
% Section style
% -------------------------------------------------
\titleformat{\section}
  {\Large\bfseries\color{titleblue}}
  {\thesection.}{0.6em}{}

\titleformat{\subsection}
  {\large\bfseries}
  {\thesubsection.}{0.5em}{}

% -------------------------------------------------
% Python code style
% -------------------------------------------------
\lstdefinestyle{pythonstyle}{
    language=Python,
    backgroundcolor=\color{codebg},
    frame=single,
    rulecolor=\color{codeframe},
    basicstyle=\ttfamily\small,
    keywordstyle=\color{keywordcolor}\bfseries,
    commentstyle=\color{commentcolor},
    stringstyle=\color{stringcolor},
    numbers=left,
    numberstyle=\tiny\color{gray},
    numbersep=8pt,
    tabsize=4,
    showstringspaces=false,
    breaklines=true,
    columns=fullflexible
}

\lstset{style=pythonstyle}

% -------------------------------------------------
% Boxes
% -------------------------------------------------
\newtcolorbox{learningbox}{
    colback=blue!4,
    colframe=titleblue,
    title=\textbf{학습 목표},
    fonttitle=\bfseries
}

\newtcolorbox{importantbox}{
    colback=yellow!8,
    colframe=orange!70!black,
    title=\textbf{핵심 포인트},
    fonttitle=\bfseries
}

\newtcolorbox{checkpointbox}{
    colback=green!5,
    colframe=green!45!black,
    title=\textbf{체크 질문},
    fonttitle=\bfseries
}

\newtcolorbox{practicebox}{
    colback=red!3,
    colframe=red!55!black,
    title=\textbf{실습},
    fonttitle=\bfseries
}

% -------------------------------------------------
% Document
% -------------------------------------------------
\begin{document}

\begin{titlepage}
    \centering
    \vspace*{3cm}

    {\Huge\bfseries Python Programming\par}
    \vspace{0.8cm}
    {\LARGE Day 1: Python 기초\par}

    \vspace{2cm}

    {\large 출력, 변수와 객체, 자료형, 연산자, 입력과 형 변환\par}

    \vfill
\end{titlepage}

\tableofcontents
\newpage

% =================================================
\section{학습 목표}
% =================================================

\begin{learningbox}
이 장을 마치면 다음 내용을 설명하고 직접 코드로 작성할 수 있다.
\begin{itemize}[leftmargin=1.5em]
    \item \texttt{print()}를 이용한 출력
    \item 변수와 객체의 관계
    \item 대입 연산자 \texttt{=}의 의미
    \item Mutable 객체와 Immutable 객체의 차이
    \item 기본 자료형 \texttt{int}, \texttt{float}, \texttt{str}, \texttt{bool}
    \item 부동소수점 오차가 발생하는 이유
    \item 산술 연산자 \texttt{/}, \texttt{//}, \texttt{\%}, \texttt{**}
    \item \texttt{input()}과 형 변환
    \item BMI 계산기 작성
\end{itemize}
\end{learningbox}

% =================================================
\section{Python과 첫 번째 프로그램}
% =================================================

Python은 문법이 간결하고 읽기 쉬운 고급 프로그래밍 언어이다.
과학 계산, 데이터 분석, 인공지능, 자동화, 웹 개발 등 다양한 분야에서 사용된다.

물리학과 공학에서는 특히 다음 라이브러리와 함께 자주 사용된다.

\begin{itemize}
    \item \texttt{NumPy}: 배열과 수치 계산
    \item \texttt{SciPy}: 과학 계산
    \item \texttt{Matplotlib}: 데이터 시각화
    \item \texttt{Pandas}: 표 형식 데이터 처리
\end{itemize}

가장 기본적인 Python 프로그램은 다음과 같다.

\begin{lstlisting}
print("Hello, World!")
\end{lstlisting}

출력 결과:

\begin{verbatim}
Hello, World!
\end{verbatim}

\begin{importantbox}
Python은 문법을 외우는 것보다, 코드가 왜 그렇게 동작하는지 이해하는 것이 중요하다.
\end{importantbox}

% =================================================
\section{출력: \texttt{print()}}
% =================================================

\texttt{print()}는 값을 화면에 출력하는 함수이다.

\begin{lstlisting}
print("Python")
print(3 + 5)
print(10 * 20)
\end{lstlisting}

출력 결과:

\begin{verbatim}
Python
8
200
\end{verbatim}

\subsection{계산과 출력의 차이}

다음 코드는 계산은 수행하지만 일반적인 Python 스크립트에서는 결과를 화면에 보여 주지 않는다.

\begin{lstlisting}
3 + 5
\end{lstlisting}

결과를 확인하려면 다음과 같이 \texttt{print()}를 사용해야 한다.

\begin{lstlisting}
print(3 + 5)
\end{lstlisting}

\begin{importantbox}
계산(computation)과 출력(output)은 서로 다른 작업이다.
\texttt{print()}는 계산 결과를 사람이 볼 수 있도록 화면에 전달한다.
\end{importantbox}

\subsection{문자열 출력}

문자열은 따옴표로 감싼다.

\begin{lstlisting}
print("Hello")
print('Physics')
\end{lstlisting}

따옴표 없이 문자를 작성하면 Python은 그것을 변수 이름으로 해석하려고 한다.

\begin{lstlisting}
print(Hello)
\end{lstlisting}

위 코드에서 \texttt{Hello}라는 이름이 정의되어 있지 않다면 오류가 발생한다.

% =================================================
\section{변수와 객체}
% =================================================

다음 코드를 보자.

\begin{lstlisting}
age = 22
\end{lstlisting}

입문 단계에서는 변수를 ``값을 저장하는 상자''라고 설명하기도 한다.
하지만 Python에서는 변수를 \textbf{객체를 가리키는 이름}으로 이해하는 편이 더 정확하다.

\[
\texttt{age} \longrightarrow 22
\]

여기서 \texttt{age}는 숫자 자체가 아니라 정수 객체 \texttt{22}를 가리키는 이름이다.

\subsection{대입 연산자 \texttt{=}}

다음 코드를 생각해 보자.

\begin{lstlisting}
age = 22
age = 23
\end{lstlisting}

두 번째 줄은 \texttt{22}를 \texttt{23}으로 수정하는 것이 아니다.

\begin{enumerate}
    \item 정수 객체 \texttt{22}가 존재한다.
    \item \texttt{age}가 \texttt{22}를 가리킨다.
    \item 정수 객체 \texttt{23}이 존재한다.
    \item \texttt{age}가 \texttt{23}을 가리키도록 바뀐다.
\end{enumerate}

즉, 바뀐 것은 객체가 아니라 \texttt{age}가 가리키는 대상이다.

\begin{importantbox}
Python의 \texttt{=}는 수학의 등호가 아니라 대입 연산자이다.
\texttt{x = 3}은 ``\(x\)와 3이 같다''가 아니라 ``\texttt{x}가 정수 객체 3을 가리키게 하라''는 뜻이다.
\end{importantbox}

\subsection{수학적 관점}

변수들의 집합을 \(V\), 객체들의 집합을 \(O\)라고 하면 프로그램의 상태를 다음과 같은 함수로 생각할 수 있다.

\[
f:V\rightarrow O
\]

예를 들어 처음에는

\[
f(\texttt{age})=22
\]

였지만, 새로운 대입 이후에는

\[
f(\texttt{age})=23
\]

이 된다.

% =================================================
\section{같은 객체를 참조하는 경우}
% =================================================

다음 코드를 보자.

\begin{lstlisting}
a = 10
b = a

a = 20

print(a)
print(b)
\end{lstlisting}

출력 결과:

\begin{verbatim}
20
10
\end{verbatim}

처음에는 \texttt{a}와 \texttt{b}가 모두 정수 객체 \texttt{10}을 가리킨다.
그러나 \texttt{a = 20}을 실행하면 \texttt{a}만 새로운 정수 객체 \texttt{20}을 가리키게 된다.

반면 다음 코드를 보자.

\begin{lstlisting}
a = [1, 2, 3]
b = a

a.append(4)

print(a)
print(b)
\end{lstlisting}

출력 결과:

\begin{verbatim}
[1, 2, 3, 4]
[1, 2, 3, 4]
\end{verbatim}

\texttt{a}와 \texttt{b}가 같은 리스트 객체를 가리키고 있고,
\texttt{append()}가 리스트 객체 자체를 수정했기 때문에 두 이름을 통해 동일한 변화가 보인다.

% =================================================
\section{Mutable과 Immutable}
% =================================================

Python 객체는 객체 자체를 수정할 수 있는지에 따라 Mutable과 Immutable로 나눌 수 있다.

\subsection{Immutable 객체}

Immutable 객체는 객체 자체를 수정할 수 없다.

대표적인 예:

\begin{itemize}
    \item \texttt{int}
    \item \texttt{float}
    \item \texttt{bool}
    \item \texttt{str}
    \item \texttt{tuple}
\end{itemize}

다음 코드는 정수 객체 \texttt{10}을 수정하지 않는다.

\begin{lstlisting}
a = 10
a = 20
\end{lstlisting}

\texttt{a}가 새로운 정수 객체 \texttt{20}을 가리키게 될 뿐이다.

\subsection{Mutable 객체}

Mutable 객체는 내부 상태를 직접 수정할 수 있다.

대표적인 예:

\begin{itemize}
    \item \texttt{list}
    \item \texttt{dict}
    \item \texttt{set}
\end{itemize}

\begin{lstlisting}
numbers = [1, 2, 3]
numbers.append(4)

print(numbers)
\end{lstlisting}

출력 결과:

\begin{verbatim}
[1, 2, 3, 4]
\end{verbatim}

\begin{checkpointbox}
다음 코드의 출력 결과를 예측해 보자.

\begin{lstlisting}
a = [1, 2, 3]
b = a
b = [7, 8, 9]

print(a)
print(b)
\end{lstlisting}

\texttt{b = [7, 8, 9]}는 기존 리스트를 수정하는 것이 아니라
\texttt{b}가 새로운 리스트 객체를 가리키게 한다.
\end{checkpointbox}

% =================================================
\section{기본 자료형}
% =================================================

Python에서 자주 사용하는 기본 자료형은 다음과 같다.

\begin{center}
\begin{tabular}{lll}
\toprule
자료형 & 의미 & 예시 \\
\midrule
\texttt{int} & 정수 & \texttt{10} \\
\texttt{float} & 부동소수점 수 & \texttt{3.14} \\
\texttt{str} & 문자열 & \texttt{"Python"} \\
\texttt{bool} & 참/거짓 & \texttt{True} \\
\bottomrule
\end{tabular}
\end{center}

자료형은 객체가 어떤 연산을 수행할 수 있는지 결정한다.

\begin{lstlisting}
a = 10
b = 3.14
c = "Physics"
d = True

print(type(a))
print(type(b))
print(type(c))
print(type(d))
\end{lstlisting}

\subsection{같은 값처럼 보여도 다른 자료형}

\begin{lstlisting}
a = 10
b = 10.0

print(type(a))
print(type(b))
\end{lstlisting}

\texttt{10}은 \texttt{int}, \texttt{10.0}은 \texttt{float}이다.
수학적으로는 같은 값이지만 컴퓨터 내부의 표현 방식은 다르다.

% =================================================
\section{부동소수점 오차}
% =================================================

다음 코드를 실행해 보자.

\begin{lstlisting}
print(0.1 + 0.2)
\end{lstlisting}

출력 결과:

\begin{verbatim}
0.30000000000000004
\end{verbatim}

이것은 Python의 버그가 아니다.

컴퓨터는 이진법을 사용하며, 십진수 \texttt{0.1}은 이진법에서 유한하게 표현되지 않는다.
따라서 컴퓨터는 표현 가능한 가장 가까운 값을 저장한다.

\subsection{수학적 모델}

실수 전체를 \(\mathbb{R}\), 컴퓨터가 표현 가능한 부동소수점 수의 집합을 \(F\)라고 하자.

그러면 저장 과정을 다음 함수로 생각할 수 있다.

\[
r:\mathbb{R}\rightarrow F
\]

컴퓨터는 실제 값 \(x\) 대신 가장 가까운 표현 가능한 값 \(r(x)\)를 저장한다.

\subsection{연산 과정의 반올림}

부동소수점 계산은 다음과 같이 생각할 수 있다.

\[
\mathrm{fl}(x+y)=r(r(x)+r(y))
\]

입력값이 저장될 때 이미 반올림이 발생하고,
연산 결과 역시 표현 가능한 가장 가까운 값으로 다시 반올림된다.

\subsection{실수 비교}

다음 두 비교는 서로 다른 결과를 낸다.

\begin{lstlisting}
print(0.1 * 10 == 1.0)
print(0.1 * 3 == 0.3)
\end{lstlisting}

출력 결과:

\begin{verbatim}
True
False
\end{verbatim}

과학 계산에서는 실수를 정확히 같은지 비교하기보다 허용 오차를 사용한다.

\[
|a-b|<\varepsilon
\]

\begin{importantbox}
부동소수점 오차는 유한한 비트로 실수를 표현할 때 발생하는 정상적인 현상이다.
\end{importantbox}

% =================================================
\section{산술 연산자}
% =================================================

\begin{center}
\begin{tabular}{cll}
\toprule
연산자 & 의미 & 예시 \\
\midrule
\texttt{+} & 덧셈 & \texttt{a + b} \\
\texttt{-} & 뺄셈 & \texttt{a - b} \\
\texttt{*} & 곱셈 & \texttt{a * b} \\
\texttt{/} & 나눗셈 & \texttt{a / b} \\
\texttt{//} & 바닥 나눗셈 & \texttt{a // b} \\
\texttt{\%} & 나머지 & \texttt{a \% b} \\
\texttt{**} & 거듭제곱 & \texttt{a ** b} \\
\bottomrule
\end{tabular}
\end{center}

\subsection{일반 나눗셈}

\begin{lstlisting}
print(10 / 3)
\end{lstlisting}

출력 결과:

\begin{verbatim}
3.3333333333333335
\end{verbatim}

\subsection{바닥 나눗셈}

\begin{lstlisting}
print(10 // 3)
print(-10 // 3)
\end{lstlisting}

출력 결과:

\begin{verbatim}
3
-4
\end{verbatim}

Python의 \texttt{//}는 단순히 소수점 이하를 버리는 연산이 아니라 바닥 함수이다.

\[
a//b=\left\lfloor\frac{a}{b}\right\rfloor
\]

\subsection{나머지}

Python의 \texttt{//}와 \texttt{\%}는 다음 관계를 만족한다.

\[
a=b(a//b)+(a\%b)
\]

예를 들어

\[
-10=3(-4)+2
\]

이므로

\[
-10//3=-4,\qquad -10\%3=2
\]

이다.

\subsection{거듭제곱}

\begin{lstlisting}
print(2 ** 3)
print(5 ** 2)
\end{lstlisting}

출력 결과:

\begin{verbatim}
8
25
\end{verbatim}

% =================================================
\section{입력: \texttt{input()}}
% =================================================

사용자로부터 값을 입력받기 위해 \texttt{input()}을 사용한다.

\begin{lstlisting}
name = input("Name: ")
print(name)
\end{lstlisting}

\texttt{input()}은 항상 문자열을 반환한다.

\begin{lstlisting}
age = input("Age: ")
print(type(age))
\end{lstlisting}

사용자가 \texttt{20}을 입력해도 자료형은 \texttt{str}이다.

\subsection{왜 문자열로 입력되는가}

키보드 입력은 숫자와 문자를 자동으로 구분하지 않는다.
사용자가 입력한 내용을 먼저 문자들의 나열로 받아들이고,
그 이후 어떤 자료형으로 사용할지는 프로그래머가 결정한다.

\subsection{형 변환}

정수로 변환하려면 \texttt{int()}를 사용한다.

\begin{lstlisting}
age = int(input("Age: "))
\end{lstlisting}

실수로 변환하려면 \texttt{float()}를 사용한다.

\begin{lstlisting}
height = float(input("Height: "))
\end{lstlisting}

수학적으로는 다음과 같은 함수 합성으로 생각할 수 있다.

\[
K
\xrightarrow{\ \texttt{input}\ }
\mathrm{String}
\xrightarrow{\ \texttt{int}\ }
\mathrm{Integer}
\]

\subsection{자료형 불일치}

다음 코드는 오류를 발생시킨다.

\begin{lstlisting}
age = input("Age: ")
print(age + 10)
\end{lstlisting}

\texttt{age}는 문자열이고 \texttt{10}은 정수이므로 두 자료형을 직접 더할 수 없다.

숫자 계산을 원한다면 다음과 같이 변환해야 한다.

\begin{lstlisting}
age = int(input("Age: "))
print(age + 10)
\end{lstlisting}

% =================================================
\section{미니 실습: BMI 계산기}
% =================================================

BMI는 다음 식으로 정의된다.

\[
\mathrm{BMI}
=
\frac{\mathrm{weight}}{\mathrm{height}^{2}}
\]

키와 몸무게를 입력받아 BMI를 계산해 보자.

\begin{lstlisting}
height = float(input("Enter your height (m): "))
weight = float(input("Enter your weight (kg): "))

bmi = weight / (height ** 2)

print(bmi)
\end{lstlisting}

예시 실행:

\begin{verbatim}
Enter your height (m): 1.80
Enter your weight (kg): 75
23.148148148148145
\end{verbatim}

\begin{practicebox}
이 프로그램은 다음 개념을 모두 사용한다.
\begin{itemize}[leftmargin=1.5em]
    \item 변수
    \item 사용자 입력
    \item 형 변환
    \item 나눗셈
    \item 거듭제곱
    \item 출력
\end{itemize}
\end{practicebox}

% =================================================
\section{핵심 요약}
% =================================================

\begin{enumerate}
    \item \texttt{print()}는 값을 화면에 출력한다.
    \item 계산과 출력은 서로 다른 작업이다.
    \item Python의 변수는 객체를 가리키는 이름이다.
    \item \texttt{=}는 수학의 등호가 아니라 대입 연산자이다.
    \item Mutable 객체는 내부 상태를 수정할 수 있다.
    \item Immutable 객체는 내부 상태를 수정할 수 없다.
    \item \texttt{int}, \texttt{float}, \texttt{str}, \texttt{bool}은 기본 자료형이다.
    \item 부동소수점 오차는 실수를 유한한 비트로 표현하기 때문에 발생한다.
    \item \texttt{//}는 바닥 나눗셈이고, \texttt{\%}는 나머지 연산이다.
    \item \texttt{input()}은 항상 문자열을 반환한다.
    \item 숫자 계산에는 \texttt{int()} 또는 \texttt{float()}를 이용한 형 변환이 필요하다.
\end{enumerate}

\end{document}
